\documentclass[10pt,letterpaper]{article}
\usepackage{cornell-notes}
\usepackage{mathtools}

\title{Chapter 1: Preliminaries}
\author{}
\date{}

\setDocTitle{Chapter 1: Preliminaries}
\setDocSubtitle{Numerical Methods}
\setDocVersion{v1.0}
\setDocStatus{Study Companion}
\setDocOwner{Numerical Methods Cornell Notes}
\setDocAuthor{}
\setDocDate{\today}
\setDocSummary{Cornell-format study and review notes for Chapter 1: Preliminaries.}

\setCornellCollection{Numerical Methods Cornell Notes}
\setCornellUnitType{Chapter}
\setCornellUnitNumber{1}
\setCornellUnitTitle{Preliminaries}
\setCornellCompanionLabel{Study and review companion}

\hypersetup{pdftitle={Chapter 1: Preliminaries},pdfsubject={Cornell notes},pdfauthor={}}

\begin{document}
\maketitle

\section*{Chapter Overview}


\section*{Learning Objectives}


\section*{Prerequisites}


\section*{Operational Workflow}


\subsection*{Formula and notation anchor}
\[
\varepsilon_{\mathrm{rel}}=\frac{|\widehat{x}-x|}{|x|},\qquad \widehat{x}=x(1+\delta)
\]
Relative error is scale-aware when the exact value is nonzero; near zero, an absolute or mixed tolerance is safer.

\begin{warnbox}{Numerical warning}
Hidden conversions, unchecked indexing, aliasing, overflow, cancellation, and abstractions that conceal allocation or copying.
\end{warnbox}

\section{1.0 Introduction}
\begin{cornell}
\noterow{What is the central idea?}{Scientific computation succeeds when mathematical intent, data representation, and software behavior are made explicit before optimization begins.}
\noterow{How should it be used?}{For Introduction, begin from explicit inputs, outputs, preconditions, and representation choices.  Specify the mathematical contract; choose representations; make dimensions, ownership, exceptional cases, and error goals explicit; then test before optimizing.}
\noterow{Cost and reliability}{Analyze arithmetic work, storage, data movement, and convergence separately.  Relevant failure pressures include hidden conversions, unchecked indexing, aliasing, overflow, cancellation, and abstractions that conceal allocation or copying.  Use scaled tolerances and report both success criteria and diagnostic evidence.}
\noterow{Implementation checklist}{\begin{itemize}
\item Validate dimensions, domains, ordering, and structural assumptions.
\item Guard small divisors, invalid states, overflow, underflow, and nonfinite values.
\item Make mutation, workspace, indexing, and termination behavior explicit.
\item Test a simple exact case, a difficult valid case, a boundary case, and an expected failure.
\end{itemize}}
\end{cornell}
\begin{summarybox}{Section summary}
Scientific computation succeeds when mathematical intent, data representation, and software behavior are made explicit before optimization begins.  Select this technique only when its assumptions match the data, and accept a result only after an independent residual, error estimate, invariant, or refinement check.
\end{summarybox}

\section{1.1 Error, Accuracy, and Stability}
\begin{cornell}
\noterow{What is the central idea?}{Accuracy measures closeness to the desired result; stability describes how an algorithm propagates perturbations; conditioning belongs to the problem itself.}
\noterow{How should it be used?}{For Error, Accuracy, and Stability, begin from explicit inputs, outputs, preconditions, and representation choices.  Specify the mathematical contract; choose representations; make dimensions, ownership, exceptional cases, and error goals explicit; then test before optimizing.}
\noterow{Quantitative lens}{Identify the governing equation, recurrence, probability law, or geometric predicate for Error, Accuracy, and Stability.  Define every variable and scale; then derive a residual, error estimate, or invariant that can be checked independently.}
\noterow{Cost and reliability}{Analyze arithmetic work, storage, data movement, and convergence separately.  Relevant failure pressures include hidden conversions, unchecked indexing, aliasing, overflow, cancellation, and abstractions that conceal allocation or copying.  Use scaled tolerances and report both success criteria and diagnostic evidence.}
\noterow{Implementation checklist}{\begin{itemize}
\item Validate dimensions, domains, ordering, and structural assumptions.
\item Guard small divisors, invalid states, overflow, underflow, and nonfinite values.
\item Make mutation, workspace, indexing, and termination behavior explicit.
\item Test a simple exact case, a difficult valid case, a boundary case, and an expected failure.
\end{itemize}}
\end{cornell}
\begin{summarybox}{Section summary}
Accuracy measures closeness to the desired result; stability describes how an algorithm propagates perturbations; conditioning belongs to the problem itself.  Select this technique only when its assumptions match the data, and accept a result only after an independent residual, error estimate, invariant, or refinement check.
\end{summarybox}

\section{1.2 C Family Syntax}
\begin{cornell}
\noterow{What is the central idea?}{Compact syntax is useful only when evaluation order, types, indexing, conversions, and ownership remain unambiguous to both reader and compiler.}
\noterow{How should it be used?}{For C Family Syntax, begin from explicit inputs, outputs, preconditions, and representation choices.  Specify the mathematical contract; choose representations; make dimensions, ownership, exceptional cases, and error goals explicit; then test before optimizing.}
\noterow{Quantitative lens}{Identify the governing equation, recurrence, probability law, or geometric predicate for C Family Syntax.  Define every variable and scale; then derive a residual, error estimate, or invariant that can be checked independently.}
\noterow{Cost and reliability}{Analyze arithmetic work, storage, data movement, and convergence separately.  Relevant failure pressures include hidden conversions, unchecked indexing, aliasing, overflow, cancellation, and abstractions that conceal allocation or copying.  Use scaled tolerances and report both success criteria and diagnostic evidence.}
\noterow{Implementation checklist}{\begin{itemize}
\item Validate dimensions, domains, ordering, and structural assumptions.
\item Guard small divisors, invalid states, overflow, underflow, and nonfinite values.
\item Make mutation, workspace, indexing, and termination behavior explicit.
\item Test a simple exact case, a difficult valid case, a boundary case, and an expected failure.
\end{itemize}}
\end{cornell}
\begin{summarybox}{Section summary}
Compact syntax is useful only when evaluation order, types, indexing, conversions, and ownership remain unambiguous to both reader and compiler.  Select this technique only when its assumptions match the data, and accept a result only after an independent residual, error estimate, invariant, or refinement check.
\end{summarybox}

\section{1.3 Objects, Classes, and Inheritance}
\begin{cornell}
\noterow{What is the central idea?}{Encapsulation can bind numerical state to valid operations, while inheritance should express a genuine substitutability relationship rather than code reuse alone.}
\noterow{How should it be used?}{For Objects, Classes, and Inheritance, begin from explicit inputs, outputs, preconditions, and representation choices.  Specify the mathematical contract; choose representations; make dimensions, ownership, exceptional cases, and error goals explicit; then test before optimizing.}
\noterow{Quantitative lens}{Identify the governing equation, recurrence, probability law, or geometric predicate for Objects, Classes, and Inheritance.  Define every variable and scale; then derive a residual, error estimate, or invariant that can be checked independently.}
\noterow{Cost and reliability}{Analyze arithmetic work, storage, data movement, and convergence separately.  Relevant failure pressures include hidden conversions, unchecked indexing, aliasing, overflow, cancellation, and abstractions that conceal allocation or copying.  Use scaled tolerances and report both success criteria and diagnostic evidence.}
\noterow{Implementation checklist}{\begin{itemize}
\item Validate dimensions, domains, ordering, and structural assumptions.
\item Guard small divisors, invalid states, overflow, underflow, and nonfinite values.
\item Make mutation, workspace, indexing, and termination behavior explicit.
\item Test a simple exact case, a difficult valid case, a boundary case, and an expected failure.
\end{itemize}}
\end{cornell}
\begin{summarybox}{Section summary}
Encapsulation can bind numerical state to valid operations, while inheritance should express a genuine substitutability relationship rather than code reuse alone.  Select this technique only when its assumptions match the data, and accept a result only after an independent residual, error estimate, invariant, or refinement check.
\end{summarybox}

\section{1.4 Vector and Matrix Objects}
\begin{cornell}
\noterow{What is the central idea?}{Vector and matrix abstractions must preserve dimensions, storage layout, indexing rules, and aliasing expectations without hiding expensive copies.}
\noterow{How should it be used?}{For Vector and Matrix Objects, begin from explicit inputs, outputs, preconditions, and representation choices.  Specify the mathematical contract; choose representations; make dimensions, ownership, exceptional cases, and error goals explicit; then test before optimizing.}
\noterow{Quantitative lens}{Identify the governing equation, recurrence, probability law, or geometric predicate for Vector and Matrix Objects.  Define every variable and scale; then derive a residual, error estimate, or invariant that can be checked independently.}
\noterow{Cost and reliability}{Analyze arithmetic work, storage, data movement, and convergence separately.  Relevant failure pressures include hidden conversions, unchecked indexing, aliasing, overflow, cancellation, and abstractions that conceal allocation or copying.  Use scaled tolerances and report both success criteria and diagnostic evidence.}
\noterow{Implementation checklist}{\begin{itemize}
\item Validate dimensions, domains, ordering, and structural assumptions.
\item Guard small divisors, invalid states, overflow, underflow, and nonfinite values.
\item Make mutation, workspace, indexing, and termination behavior explicit.
\item Test a simple exact case, a difficult valid case, a boundary case, and an expected failure.
\end{itemize}}
\end{cornell}
\begin{summarybox}{Section summary}
Vector and matrix abstractions must preserve dimensions, storage layout, indexing rules, and aliasing expectations without hiding expensive copies.  Select this technique only when its assumptions match the data, and accept a result only after an independent residual, error estimate, invariant, or refinement check.
\end{summarybox}

\section{1.5 Some Further Conventions and Capabilities}
\begin{cornell}
\noterow{What is the central idea?}{Consistent naming, types, exceptions, templates, and helper conventions reduce accidental complexity in later numerical implementations.}
\noterow{How should it be used?}{For Some Further Conventions and Capabilities, begin from explicit inputs, outputs, preconditions, and representation choices.  Specify the mathematical contract; choose representations; make dimensions, ownership, exceptional cases, and error goals explicit; then test before optimizing.}
\noterow{Quantitative lens}{Identify the governing equation, recurrence, probability law, or geometric predicate for Some Further Conventions and Capabilities.  Define every variable and scale; then derive a residual, error estimate, or invariant that can be checked independently.}
\noterow{Cost and reliability}{Analyze arithmetic work, storage, data movement, and convergence separately.  Relevant failure pressures include hidden conversions, unchecked indexing, aliasing, overflow, cancellation, and abstractions that conceal allocation or copying.  Use scaled tolerances and report both success criteria and diagnostic evidence.}
\noterow{Implementation checklist}{\begin{itemize}
\item Validate dimensions, domains, ordering, and structural assumptions.
\item Guard small divisors, invalid states, overflow, underflow, and nonfinite values.
\item Make mutation, workspace, indexing, and termination behavior explicit.
\item Test a simple exact case, a difficult valid case, a boundary case, and an expected failure.
\end{itemize}}
\end{cornell}
\begin{summarybox}{Section summary}
Consistent naming, types, exceptions, templates, and helper conventions reduce accidental complexity in later numerical implementations.  Select this technique only when its assumptions match the data, and accept a result only after an independent residual, error estimate, invariant, or refinement check.
\end{summarybox}

\section*{Method comparison}
\begin{center}
\small
\begin{tabularx}{\linewidth}{>{\bfseries\raggedright\arraybackslash}p{0.22\linewidth}XX}
\toprule
Method or lens & Best use & Main caution \\
\midrule
Absolute error & Values near zero & Depends on physical scale \\
Relative error & Nonzero scaled values & Unstable as the reference approaches zero \\
Backward error & Algorithm assessment & Must be interpreted through conditioning \\
\bottomrule
\end{tabularx}
\end{center}

\section*{Worked example}
\begin{exambox}{Independent worked example}
For $x=10/3$ and $\widehat{x}=3.333$, the absolute error is about $3.33\times10^{-4}$ and the relative error is about $10^{-4}$.  A tolerance test should combine absolute and relative scales: $|e|\leq a+r|x|$.  The check avoids division by a nearly zero reference; finite precision can still round the computed error itself.
\end{exambox}

\section*{Chapter synthesis}
\begin{summarybox}{Chapter synthesis}
The chapter's methods should be viewed as a decision system rather than a list of routines.  Begin with the mathematical problem and its conditioning, exploit structure, select an algorithm with compatible stability and cost, and verify the computed answer with evidence that is independent of the internal iteration.  The recurring implementation discipline is: specify the mathematical contract; choose representations; make dimensions, ownership, exceptional cases, and error goals explicit; then test before optimizing.
\end{summarybox}

\subsection*{Common mistakes and numerical hazards}
\begin{warnbox}{Numerical warning}[title={Common mistakes and numerical hazards}]
Hidden conversions, unchecked indexing, aliasing, overflow, cancellation, and abstractions that conceal allocation or copying.  A second recurring mistake is reporting only a computed value: always retain tolerances, iteration counts, residuals, refinement behavior, and relevant structural checks.
\end{warnbox}

\subsection*{Retrieval practice}
\textbf{Questions}
\begin{enumerate}
\item What problem does Error, Accuracy, and Stability address, and which assumption is easiest to violate?
\item What problem does C Family Syntax address, and which assumption is easiest to violate?
\item What problem does Objects, Classes, and Inheritance address, and which assumption is easiest to violate?
\item What problem does Vector and Matrix Objects address, and which assumption is easiest to violate?
\item What problem does Some Further Conventions and Capabilities address, and which assumption is easiest to violate?
\item How do conditioning, stability, and the final residual play different roles in this chapter?
\end{enumerate}

\textbf{Brief answers}
\begin{enumerate}
\item Accuracy measures closeness to the desired result; stability describes how an algorithm propagates perturbations; conditioning belongs to the problem itself. Recheck the section's domain, structure, scaling, and failure diagnostics.
\item Compact syntax is useful only when evaluation order, types, indexing, conversions, and ownership remain unambiguous to both reader and compiler. Recheck the section's domain, structure, scaling, and failure diagnostics.
\item Encapsulation can bind numerical state to valid operations, while inheritance should express a genuine substitutability relationship rather than code reuse alone. Recheck the section's domain, structure, scaling, and failure diagnostics.
\item Vector and matrix abstractions must preserve dimensions, storage layout, indexing rules, and aliasing expectations without hiding expensive copies. Recheck the section's domain, structure, scaling, and failure diagnostics.
\item Consistent naming, types, exceptions, templates, and helper conventions reduce accidental complexity in later numerical implementations. Recheck the section's domain, structure, scaling, and failure diagnostics.
\item Conditioning describes sensitivity of the mathematical problem, stability describes error propagation by the algorithm, and the residual measures satisfaction of the computed equations; none is a universal substitute for the others.
\end{enumerate}

\subsection*{Mastery checklist}
\begin{itemize}
\item[$\square$] I can explain and select Error, Accuracy, and Stability without relying on a memorized code listing.
\item[$\square$] I can explain and select C Family Syntax without relying on a memorized code listing.
\item[$\square$] I can explain and select Objects, Classes, and Inheritance without relying on a memorized code listing.
\item[$\square$] I can explain and select Vector and Matrix Objects without relying on a memorized code listing.
\item[$\square$] I can explain and select Some Further Conventions and Capabilities without relying on a memorized code listing.
\item[$\square$] I can state a scaled stopping rule and an independent verification check.
\item[$\square$] I can identify at least one conditioning risk and one algorithmic stability risk.
\item[$\square$] I can estimate time, storage, and data-movement costs at the appropriate level.
\end{itemize}

\end{document}


